\section{Lecture 27}
\subsection{Main notes}
Faraday's law
\[
\epsilon = - \dv{\Phi_B}{t} = -L \dv{I}{t}
  \]
\[
L = \mu_0 NA/L
  \]
\[
C = \frac{Q}{V}
  \]

  Combining L and C gives us an LC circuit, which can be modeled as a harmonic ossilator
\[
w = \frac{1}{\sqrt{LC}}
  \]
  If one combined this with an AC source, LCR circuit.

  We find thatin this case, as R approches $X_R = R$. We can find $V_{max} = I_{max}R$

  As C approches $X_c = \frac{1}{wc}$. We get $V_{max} = I_{max}X_C$

As L approches $X_L = W$ We can find $V_{max} = I_{max}X_L$

\[
Z = \sqrt{R^2 + (X_L^2 + X_C^2)}
  \]

\[
V = IZ
\]

\[
P=iv
\]

Already doing electromagnetic waves.

For polarization, it is the difference between the angle and the polarized axis. Make sure to check if the wave is polarized or not.
\begin{itemize}
	\item
\end{itemize}


\subsection{Connections}


\subsection{Summary}


\subsection{Questions}
